\section{帰無仮説検定}
\label{lesson18}
\subsection{授業内容}
\subsubsection{科目の中でのこのコマの位置づけ}
\begin{tabular}[t]{cccccc}
    記述統計[4] & 確率[5] & モデル[9] & 推測・推定[2] & 意思決定[2] & 実践[2]
\end{tabular}
\subsubsection{概要}
心理学の研究において最もよく使われる意思決定法である，帰無仮説検定の具体例について学ぶ。帰無仮説検定は背理法をもちいて効果の有無を判断する方法であり，結果の報告に関する表記法も含めてその作法を学ぶ必要がある。注意するべきは，帰無仮説検定の技術は誤用されたりや悪用されたりすることが多い方法であり，前提や仮定をしっかり守って正しく運用することを常に心がけねばならない点である。ここでは最も単純な相関係数の帰無仮説検定に基づいて考える。また，検定は確率的判断の手続きであるが故に，判断の誤り方が2通り考えられる。ここで検定力の考え方が導かれ，どのような状況のもとで帰無仮説検定が行われているかを理解する必要がある。
\subsubsection{コマ主題細目}
\begin{description}

\item[1つの平均値の検定]ここでは母分散がわかっている場合の，標本平均の検定を考える。標準化の手続きと，帰無仮説検定の考え方を一通り復習した上で，実際どのように検定が行われるかの実践例である。
\begin{flushright}
    $\to$\citet{川端201408}のP.80--89，\citet{Yamada2004}のP.126--127.
\end{flushright}

\item[1つの平均値の検定(母分散未知)]先ほどの例では，検定すると言いながら母分散がわかっているという特殊な状況であった。普通は母分散がわかっていることはなく，これも不偏分散を使って推定しなければならないが，このような場合は標準正規分布ではなく，$t$分布と呼ばれる統計量が用いられる。$t$分布の導出は行わないが，不偏分散を使った分布であること，サンプルサイズによる大きさの違いを調整する，自由度というパラメータがあることに注意する。
\begin{flushright}
    $\to$\citet{川端201408}のP.93--96.
\end{flushright}

\item[結果の報告]相関係数の有意性検定について解説する。別名「無相関検定」ともよばれているように，帰無仮説が想定する世界は母相関が0の状態である。ここから得られた標本が作る標本相関の分布は$t$分布に従うことがわかっており，これを使って検定を行うことができる。計算のプロセスを簡単にフォローするとともに，結果の報告における表記法について理解する。
\begin{flushright}
    $\to$\citet{Yamada2004}のP.132--133,\citet{山内201001}のP.216--220.
\end{flushright}

\item[有意はえらい？]相関係数の帰無仮説検定においては，顕著な違いが示されること＝相関係数が大きいこと，を意味するわけではない。例えばサンプルサイズが大きくなれば，非常に小さな相関係数であっても統計的に有意であると判断される。有意かどうかにこだわるのではなく，その大きさに基づいて判断することの重要性を学び，標準化された効果の大きさを表現する効果量について理解する。
\begin{flushright}
    $\to$\citet{Yamada2004}のP.142--143.効果量については\citet{大久保201201}のP.43--46.
\end{flushright}


\item[検定における2種類の間違い]ここで検定における2種類の間違い方，すなわちType I ErrorとType II Errorについて解説する。モデルの定量的な比較でなく，仮定に基づく質的な確率的な判断であり，その組み合わせから4つの結果が起こり得る。ここでNHSTでは何をコントロールし，何を目的としているかを改めて確認すること，また検定を繰り返すことでType I Errorがインフレーションを起こしうることを理解することが必要である。
\begin{flushright}
    $\to$ \citet{Yamada2004}のP.120--121，検定の多重性については同書P.158--161.\citet{川端201408}の P.90--93.も
\end{flushright}

\end{description}
\subsubsection{キーワード}
\begin{itemize}
\item 一つの平均値の検定
\item $t$分布
\item 無相関検定
\item 信頼区間
\item 効果量
\end{itemize}
\subsection{授業情報}
\paragraph{コマの展開方法}
適宜投影資料を用いる
\paragraph{標準シラバスにおける位置づけ}
科目番号4；心理学研究法；2.データを用いた実証的な思考方法；C データの統計的記述\\
科目番号5；心理学統計法；1.心理学で用いられる統計手法；F 推測統計:その考え方
\subsubsection{予習・復習課題}
\paragraph{予習}
具体的な帰無仮説検定の手続きが進むので，帰無仮説，対立仮説，有意水準といった一般的な用語について復習しておく。
\paragraph{復習}
相関係数の大きさ，標本の大きさを任意の値にし，無相関検定を実際にやってみると良い。様々な数値を入れることで，検定結果がどのように変動するかを理解することは帰無仮説検定の本質を理解するのに役立つ。仮想的な数字では実感が得られにくい場合，身の回りの相関がありそうなデータを具体例として検定し，何が言えるのかを考えてみると良い。